
% 来源：EPUB part056.xhtml
\section{估值实际上是“功夫”}

\subsection{定性比定量分析重要}

有很多问题是关于估值的、所以简单谈谈自己的想法。

我个人觉得如果需要计算器按半天才能算出来那么一点利润的投资还是不投的好。

我认为估值就是个毛估估的东西、如果要用到计算器才能算出来的便宜就不够便宜了。

好像芒格也说过，从来没见巴菲特按着计算器去估值一家企业，我好像也没真正用过计算器做估值。

我总是认为大致的估值主要用于判断下行的空间，定性的分析才是真正利润的来源，这也可能是价值投资里最难的东西。

一般而言，赚到几十倍甚至更多的股票绝不是靠估值估出来的，不然没道理投资人一开始不全盘压上。（当时我要知道网易会涨160倍，我还不把他全买下来？）

正是由于定性分析有很多不确定性，所以多数情况下人们往往即使看好也不敢下大注，或就算下大注也不敢全力以赴。

当然，确实也有一些按按计算器就觉得很便宜的时候，比如巴菲特买的中石油、我买的万科。但这种情况往往是一些特例。

我觉得估值在决定卖股票时也是毛估估的，比如当万科涨到3500-3600多亿时，我曾强烈建议我一个朋友应该卖掉，因为我觉得万科要持续赚300多亿是件很困难的事，你花3000多亿去赚不到100亿的年利润实在不是个好主意。如果你真的喜欢万科，可以等一等，将来可能会有机会用同样的钱买更多的股份回来。说实话，我挺喜欢万科的，我想有一天我可能会再投万科的。下次投的时候争取多买些，拿的时间长一点。

巴菲特确实说过伟大的公司和生意是不需要卖的，可他老人家到现在为止没卖过的公司也是极少的。另外，我觉得巴菲特说这话的潜台词是其实市场往往不会给伟大的公司一个疯狂的价钱，如果你仅仅是因为有一点点高估就卖出的话，可能会失去买回来的机会。而且，在美国，投资交的是利得税，不卖不算获利，一卖就可能要交很高的税，不合算。以上纯属个人猜测。（2010-04-24）

其实现金流折现法（Discounted Cash Flow, DCF）就是毛估估的，所以是一个东西。

我非常同意DCF（生命周期的总现金流折现）是唯一合乎逻辑的估值方法的说法，其实这就是“买股票就是买公司”的意思，不过是量化了。（2010-04-27）

我觉得“模糊的正确”指的就是定性的重要性，或者甚至也可以说是要先做对的事情。（2018-05-28）

我所谓量的说法都是建立在定性分析上的。从10年或以上的角度看，定量分析很荒唐。（2019-08-21）

定性分析要更准确。比如人们可能很容易知道茅台10年后大概能赚多少钱，但无法准确知道10年后的某一年能赚多少，也没必要知道。对苹果这类公司则更是如此。（2019-08-23）

\BookQuestion{网友：有人说估值是门艺术，我觉得估值是科学加艺术，挺难的。做投资的人都要解决这个课题。}

我也不知道估值属于科学还是艺术，我还一直以为科学就是艺术，艺术也挺科学的呢。我总觉得估值实际上是“功夫”。《功夫熊猫》电影里最大的缺憾是功夫没有体现时间性。实际上中文里功夫本身就是时间的意思。（2011-01-05）

\BookQuestion{网友：估值就是时间，我不太明白，可能段先生从《功夫熊猫》里悟出的东西，凭我的水平还不能理解。}

大家别把《功夫熊猫》太当回事了，那是娱乐产品。

我没说估值就是时间啊。我的言下之意是：

1. 要学会估值是要花很长时间的，绝不是很多人以为的不知道哪里有一个神秘的公式，每当要估值时用公式算一下就可以了。这个很长的时间大概不是用小时或天来计算的，用很多年作为单位比较合适。

2. 对某个企业的估值经常也是要花很长的时间的。

我对估值的定义基本上就是对企业的了解，只有当我觉得很了解很了解一个企业的时候，我才能对企业有一个大概的估值，这往往需要很长的时间。

当然，过去对企业的了解的积累也是非常有用的，这就是有时候好的投资决策看起来并没有花很多时间来决定，实际上却是很多年了解的积累的结果。比如，巴菲特决定投高盛可能只花了20分钟，但那其实是他50年理解的积累。我当年“敢”重仓网易也不是看起来那么容易的。我错失苹果也同样是因为积累不够，找不到很了解的感觉。

总而言之，估值是需要“功夫”的。（2011-01-08）

% 来源：EPUB part057.xhtml
\subsection{看财报主要用于排除公司}

\BookQuestion{网友：您会把财报看得很细吗（除了三大表外）？会经常参加股东大会吗？}

我不是很经常看财报（所以说没有巴菲特用功）。比较在意的数字是几个：负债、净现金、现金流、开销合理性、真实利润、扣除商誉的净资产。其他好像没了。

不知道三大表是什么。我看财报主要用于排除公司，也就是说如果看完财报就不喜欢或看不懂的话，就不看了。决定投进去的原因往往是其他的因素，UHAL是个例外。

也有没看过财报就投的，比如伯克希尔A。去年买了些伯克希尔A并且做成了证书，就为了让巴菲特签个名，呵呵。结果到现在为止，已经赚了几顿饭钱了。很少参加股东大会，但基本每年都去奥马哈。（去年由于流感没去。）（2010-03-08）

巴菲特曾经向股东推荐了几本书，有一本是《杰克·韦尔奇自传》，你去看那本书，你会发现韦尔奇对企业的文化问题是很在意的，所以你可以想象巴菲特对此也很在乎。他不是像有的人说的只看财报。只看财报只会看到一个公司的历史。我看财报不读得那么细，但我找专业人士看，别人看完以后给我一个结论，对我来讲就OK了。但是我在乎利润、成本这些数据到底是由哪些东西组成的，要知道它真实反映的东西是什么。而且要连续几个季度甚至几年来看，跟踪一家公司久了，你就知道它是在说谎还是说真话。

好多公司看起来赚很多钱，现金流却一直在减少，那就有危险了。其实这些巴菲特早说过，人们都知道，但是投资的时候就会糊涂。很多人管他们自己叫投资，我却说他们只是for fun，他们很在乎别人满不满意（这家公司），真正的投资者绝对是“目中无人”的，脑子里盯的就是这个企业，他不看周围有没有人买，他最好希望别人都不买。同样，如果我做一个上市公司，我也不理（华尔街）他们，该干嘛干嘛，股价高低跟我没有关系的。

所以我买公司的时候，有一个很大的鉴别因素就是，这家公司的行为跟华尔街对它的影响有多大的关联度。如果关联度越大，我买它的机会就越低。华尔街没什么错，华尔街永远是对的，它永远代表不同人的想法。但是你要自己知道自己在干什么，如果自己没了主见，要听华尔街的，你就乱了。（2013-09-02）

\BookQuestion{网友：您有没有鉴别财务造假化繁为简的招数？}

我投资是不太看财报的（早期看过一点点，不觉得有用）。一般来说，不了解的公司不要碰，了解的公司不太可能不知道他们是不是有意造假的。如果你连业务的造假都看不到，那财报造假就更难看到了。（2019-09-19）

% 来源：EPUB part058.xhtml
\subsection{不产生现金流的净资产没有价值}

计算价值只和未来总的现金流折现有关。其实净资产只是产生未来现金流的因素之一，所以我编了个“有效净资产”的名词。也就是说，不能产生现金流的净资产其实没有价值（有时还可能是负价值）。（2010-05-08）

不赚钱的净资产有时候就是个累赘。比如在渺无人烟的地方建个酒店，花了一个亿，现在每年亏500万，重置成本还是一个亿，现在5000万想卖，这里谁要？（2010-04-25）

\BookQuestion{网友：轻资产的公司是不是给高点溢价？像贵州茅台的无形资产是很值钱的。}

无形资产也是可以折现的，所以我不建议单独考虑无形资产，因为公司的获利能力里已经包含了。（2010-09-23）

\BookQuestion{网友：现金流应关注哪些方面，才能看出公司在作假？}

假的真不了，时间会告诉你，但好像没有公式，否则就没人能作假了。不过，了解其产品、市场、竞争对手以及熟悉财务报表可能都能帮你发现可能的作假。（2010-03-14）

\BookQuestion{网友：我看了一家公司过去六七年的报表，现金流的总和是利润总和的一半左右。如果一家公司持续这样，您会看它的利润还是现金流？}

如果现金流少于利润的话就要小心了，要看它的钱花在哪里，是否确实对未来的发展有好处，不然会不见的。长期现金流少于利润不是一件好事，要仔细看财报。现金流多于利润往往是由于折旧或预收等原因。不管是什么原因，最好能明白。

\BookQuestion{网友：我能看到大概是存货和应收账款变多了。以后看到现金流长期少于利润的我直接把它过滤掉，不用小心了，直接省心。您常说的现金流的确是财务上的那个现金流吧？}

还有别的意思吗？一般财务不理解的地方是，时间足够长的情况下，现金流实际上就是利润（或叫基本等于），你想想看什么道理。我有时会把两个东西混一起说是因为我总是假设很长的时间。如果只看一年或几年，它的差别可以很大，有差别时往往要小心，尤其是现金流老是远小于利润时。但是，有些公司在投放期的利润往往是大于现金流的，比如淘宝头几年，现金流有可能是负的（花的钱比进来的钱多，不算发行股票的话），现在情况怎么样不知道。（2010-05-27）

\BookQuestion{网友：您现在是否依然坚持这个观点？您怎么看拼多多依然没有盈利这个问题？}

“不赚钱的生意多少营业额都是没用的。”这个观点有啥问题吗？如果你打算买下整个公司你就会明白的。当然，这里指的是公司的整个生命周期，和某个公司没有关系。比如，如果拼多多最后无法盈利，那这个公司的价值最后就会是零。（2020-10-17）

黄金不产生现金流，不保值。（2020-05-17）

\BookQuestion{网友：黄金是不是能跑赢通胀？}

我以前也这么想过，后来发现黄金是跑不过通胀的（黄金的购买力也是会下降的）。不过，恶性通胀发生的时候，黄金应该有一定的保值作用吧。（2022-01-25）

比特币不产生现金流，无法评估，属于不懂不碰的范畴。同理，我也不会投资名画、邮票、古董、郁金香……自己用的不算投资，虽然也可能会赚钱。（2022-05-02）

前几天还跟很早就买了不少比特币的一个朋友聊过几句。他的观点是比特币就像黄金，和市面上的说法差不多。我似懂非懂。黄金是稀缺的，比特币也是稀缺的，但搞出一个像比特币一样稀缺的东西很容易啊，搞出一个黄金属性的东西很难，毕竟铜汞银铂金，金是最稳定的金属哈。（2020-11-23）

净资产收益率（Return on Equity, ROE）是个很好的指标，可以用来排除那些你不喜欢的公司，但不能作为“核心指标”去决定你是否要投这家公司。这个指标如果好的话，你首先要看的是债务；如果没有债务问题则马上要看文化，看这种获利是否可以持续。最后还是要回到对的生意，对的人，对的价格。（2015-02-12）

在没有贷款的前提下，净资产收益率高的公司当然是不错的，说明公司的盈利能力强。净资产收益率低的公司一般都不太好，但成长型公司的初期可能净资产收益率会很低。关键还是看你能不能看懂公司的未来（现金流）。（2011-02-13）

\BookQuestion{网友：您持有的苹果、茅台、腾讯都是净利润率特别高的公司，高净利率的公司是否就比低净利率的公司商业模式好？}

好像我喜欢的公司都有高毛利的特征，但没归纳过。

高毛利肯定不是充分条件，比如餐饮业都是高毛利，但投资价值不高。（2019-03-19）

\BookQuestion{网友：我们公司最近要融资，投行建议我们使用的估值依据是息税折摊前利润（EBITDA）。记得芒格很干脆说过EBITDA就是狗屎。EBITDA的参考价值到底在哪儿呢？}

芒格都说EBITDA是狗屎咯，狗屎的参考价值是什么？（2015-12-02）

“如果我们把世界上所有谈论EBITDA的人和世界上所有没有谈论EBITDA的人放在一起，第一组人中会有更多的欺诈，并且从百分比上看，两组人会有非常大的差距。我是说，问题还只是刚刚开始。”（巴菲特）（2024-02-08）

很赚钱的公司是没道理用的，除非他们用错了会计师（事务所）。（2024-02-09）

千万不要被P/E（Price- to- Earnings Ratio，市盈率）误导，因为P/E是个历史数据，未必能说明未来。

我说P/E一般指的是相对于未来长期实际利润的P/E，不是一般财报上的那个P/E。

你愿意给多少P/E完全取决于你自己的能力或者说你自己的资金的机会成本，其实和市场无关。（2011-10-01）

多少倍P/E和平均长期利息有关，12倍长期P/E应该可以了，我一般也用这个数左右。（2011-01-05）

\BookQuestion{网友：曾经在A股满仓民生银行的我，现在看着满仓的茅台，感谢大道。}

让我想起很多年前投了很多钱在民生银行的一个熟人，那时要都换成茅台现在应该会舒服很多哈。（2023-03-15）

\BookQuestion{网友：投资方面的认知，史当时确实要弱一些。记得他说过一句让我印象很深的话：“15倍以上市盈率的，我不买。”}

我买茅台的时候好像市盈率也不高吧？我不太记得，只是关注商业模式了。市盈率太高不买多数情况下是对的，我也不喜欢买市盈率高的公司，除非看到未来了。（2023-03-15）

% 来源：EPUB part059.xhtml
\subsection{是否分红回购不是回报率的关键}

其实是否回购或分红都不是回报率的关键，但真正能赚到钱的好公司最后都一定会通过某个或某几个途径将利润给到股东的，除非公司能像巴菲特那样用其多余的钱持续去做回报高于社会平均水平的投资……总有一天巴菲特的公司也会有用不完的钱的。（2015-11-16）

国内上市公司的一个怪现象就是又借钱又分红，因为不分红就不能融资，这是很多年前我问王石时知道的。这种类型的分红是监管的问题，有点怪异。（2013-03-08）

\BookQuestion{网友：这些年巴菲特对IBM谈得也挺多。现在的IBM，除了价格比4年前还低，股票总数也已经减少了将近25\%！}

股票减少的同时债务增加了多少？靠债务回购是件非常不舒服的事情。（2015-04-08）

\BookQuestion{网友：巴菲特特别喜欢公司回购股票，回购会提升股票的内在价值吗？}

回购其实也不提高整个公司的内在价值，但提高个股的内在价值，因为股票少了。（2011-03-03）

别的公司回购的真实理由可能有无穷个，但巴菲特回购的理由大概就只有一个，那就是他现在暂时很难找到比伯克希尔更便宜的股票，也就是说他认为伯克希尔确实很便宜了。我理解的巴菲特对很便宜的定义大概是安全的年复合回报率在10\%或以上的意思。（2011-09-26）

\BookQuestion{网友：“安全的年复合回报率在10\%或以上”，这个数据是怎么得来的？}

这个数字是猜的。巴菲特一般安全回报8\%或以上的可转换债券就愿意投了。（2011-09-29）

% 来源：EPUB part060.xhtml
\subsection{股票由每个买家自己“定价”}

股票是由每个买家自己“定价”的，到你自己觉得便宜的时候才可以买，实际上和市场（别人）无关。啥时你能看懂这句话，你的股票生涯基本上就很有机会持续赚钱了。如果看不懂其实也没关系，因为大概85\%的人是永远看不懂这句话的，这也是这些人早晚会亏在股市上的根本原因。非常有趣的是，我发现其实大部分从事“投资”行业的“专业”人士其实也不是真的很明白这句话。（2012-04-14）

我其实并不懂通常意义上的估值，就是所谓股价应该是多少的那种估值。我一般只是想象如果某个公司是非上市企业，用目前的市值去拥有这家公司和其他的机会比较，哪个在未来10年或20年得到的可能回报更高（这里的回报是指公司的盈利而不是股价的涨幅）。（2019-03-13）

\BookQuestion{网友：我看见一只股票，它以前的十年净收益从来没有超过1.5亿，但它的市值有100亿。每个人都知道它的前景是一般般罢了。为什么那么多人去抢它，难道每个人都不懂吗？}

呵呵，这有啥不明白的，再不明白去趟澳门就什么都明白了。（2010-05-14）

我个人理解逆向思维指的是从不同的角度看问题。在投资里，逆向思维大概就是能看到大家看不到的价值吧，不然大家也不会用很便宜的价钱卖给你。（2012-05-15）

\BookQuestion{网友：我去过几次股东大会，分别问过董事长、总经理、董秘，他们都说自己公司的价格被低估了。}

大概所有公司的人都会对外说自己的股票被低估了吧？问完之后你感觉好点儿了吗？（2012-02-13）

\BookQuestion{网友：大股东的回购能视为股票被明显低估吗？}

回购最多只能说明大股东认为其股票被低估了。很多时候公司的回购只是为了给市场一个姿态，实际上并不表示什么。大概只有一种情况能表示公司确实认为其股票真的被低估了，那就是完全私有化。即便如此，也不代表这个股票真的被低估了。（2012-02-11）

% 来源：EPUB part061.xhtml
\subsection{案例}

\BookMinorHeading{UHAL}

\BookQuestion{网友：香港现在有些烟蒂资产。营收100多亿，经营现金流为5亿以上，利润这几年算正负零，有80亿的净资产，但市值才10个亿。这种公司怎么看？记得您买过UHAL？}

我现在不再买这种生意了。（2020-10-23）

我的过滤器里有商业模式和企业文化两项，过不去的我就大概看不到了。UHAL某种角度属于烟蒂。（2020-10-25）

\BookQuestion{网友：您有投资过隐性资产没有在账面上体现的上市公司吗？}

有。雅虎就是，当年UHAL也是。（2010-03-27）

\BookQuestion{网友：请问怎样才能找到这种好事呢？}

我是朋友告诉我的。我告诉过很多朋友，好像没人把这个当好事。（2010-04-25）

\BookQuestion{网友：我搞不懂为什么报表会不把全部的资产计算在内。怎么分辨这种情况？}

有时候会有，比如当年我买UHAL时就是这种情况。UHAL拥有的地产很值钱，但报表没法反映出来。巴菲特买的BNI我觉得也有这种资产，但比例好像没那么高。（2010-05-11）

\BookQuestion{网友：您找人核算了一遍UHAL的资产吗？想了解一下您的算法。}

我的算法很简单，就是朋友告诉我UHAL有50美元以上的净资产，我不太相信，因为当时股价在5美元左右，然后找不同的人查了一下，自己也看了看，确实是有50美元左右，就买了。这种机会就是巴菲特讲的天上掉馅饼的事吧？最关键的是这时候你的水桶要够大，不然也接不到啥。（2010-03-09）

\BookQuestion{网友：为什么这么便宜？}

他们当时就是现金出问题了。（2010-04-17）

\BookQuestion{网友：您真找会计师看过报表？}

是真找人看了。不过最关键的是一个朋友告诉我这么大个馅饼在那儿，我不太信，所以才找人看其报表，主要是看它的资产是不是真有那么多。看的结果和朋友告诉我的一样。后来再从侧面去了解了一下这个公司，也大致明白为什么市场不喜欢，就毛估估下手了。遗憾的是正好买在了底部，所以买得不够多。

\BookQuestion{网友：“UHAL许多人唯恐避之不及，也就是这时候，段永平请会计师审查财务报表，确定公司业绩没有造假”——光看财报无法确定公司业绩没有造假吧？而公司的原始凭证、业务往来之类如果它不配合也很难进行审计，对吧？}

找人看主要是看资产，不是看真假。美国这点好，假的概率小。UHAL的资产情况很复杂，花了很久（几个礼拜）才全面搞明白。UHAL现在好像又快100美元了。（2012-02-02）

% 来源：EPUB part062.xhtml
\BookMinorHeading{万科}

我现在更关注商业模式和企业文化。如果时光能够倒流，我大概不会买创维或者万科，我觉得它们的商业模式都不够好。（2021-12-30）

\BookQuestion{网友：万科109亿股本，市值870亿，存货900亿，且账上有230亿的资金，您有什么看法？}

之前的经营和负债情况如何？（2010-04-23）

\BookQuestion{网友：万科的负债，2012年长期 + 短期的负债就达到了460亿。不知道净负债率23.5\%是怎么算出来的？}

你要算他们的有息负债。预收款也是负债，但意义差很多。（2013-03-06）

\BookQuestion{网友：不只是万科有负债，每个公司都有负债，只要是负债不是很离谱就行了。}

每个公司都有负债吗？每个公司都应该负债吗？每个房地产公司都应该负债吗？我记得美国有一家房地产公司，叫NVR，曾经因为太多负债而差一点破产，股票掉到过0.2美元/股，后来侥幸活过来以后就再也不借钱了，几十年以后的现在股价大概在700美元/股上下。你再看看美国房地产行业的其他公司。如果万科没负债，我会给他们加很多分，但他们不懂，他们想快点发展嘛。（2011-09-27）

\BookQuestion{网友：请问是不是关注有息金融负债就行了，其他负债就不用管了。我看到万科的有息负债只有200多亿，而它却有400多亿现金，我以为很安全。}

还没认真看过财报。如果现金超过负债，那就应该很安全。我对大多数公司的担心是很多公司为了发展快些而去负债，结果某一天因为负债而陷入麻烦。万科相对而言还是比较健康的，但似乎对营业额的增长也很有追求。我一般喜欢那些追求好产品，不太提营业额目标的公司，比如苹果。

\BookQuestion{网友：现金400多亿，搞不懂有息负债200多个亿怎么来的？}

有息负债大概就是银行借贷或债券类的东西，到期要还的。我说的少借债指的就是这类债。会计上的债务还包括预收类的款项，这类债务一般不会给企业带来风险。（2011-09-28）

\BookQuestion{网友：现金超过负债是不是有点怪啊。自己有钱放着不花，借别人的钱给别人利息？}

很多很大的公司（比如跨国公司）由于周转等各种原因，有时候会有短期贷款的需求。（2011-09-30）

\BookQuestion{网友：感觉A股很少公司不负债的，不知道美国怎么样？}

没统计过。反正一般负债多的公司我也不买。（2010-04-24）

\BookQuestion{网友：我个人偏爱低负债、零负债的公司。还要看看资本开支，有不少公司是资本黑洞，利润只是纸上富贵。有些公司多年以来的全部利润再投入都不够，还要增发或举债，我会避开此类公司，哪怕利润表很漂亮。}

我也尽量避开。（2010-04-04）

\BookQuestion{网友：您认为一线城市的房地产符合巴菲特所说的“长坡厚雪”吗？}

我不懂房地产，对这个行业普遍高举债感到困惑。（2019-04-22）

\BookQuestion{网友：像万科等地产公司的经营性现金流经常为负，你觉得这样危险吗？}

对有些公司而言总有一天会的。如果整个行业都如此，则对有些公司而言后果很严重。（2010-09-03）

\BookQuestion{网友：你当初买万科时是怎样估值的？}

我的所谓估值都是毛估估的。我买万科时万科市值才100多个亿，我认为无论如何也不止，所以就随便给了个500亿。呵呵，当时唯一确定的就是100多个亿有点太便宜了，如果有人把万科100多个亿卖给我，我会很乐意买，所以买一部分也是很乐意的。（2010-03-01）

\BookQuestion{网友：万科总股本110亿，今日股价7.31元，总市值804亿，每股净资产3.55元，净资产总计390.5亿，资产负债67\%。}

我确实认为万科是一家相对不错的公司，但我也不知道什么价该买。去年我曾经设了6块钱的买价，结果他掉到6块多就又跑回去了，呵呵。即使是5年10年以后，人们也还是要买房子的，他们会买谁造的房子呢？

\BookQuestion{网友：6块估价怎么来的？}

依据记得6块钱时市值大概800亿上下，加上贷款1000亿出头，100亿左右年利润，未来经营走下坡的机会不大，这个价钱我觉得不贵。（2011-03-08）

\BookQuestion{网友：市值为什么还要加上贷款呀？想不明白。}

假设你按市值买下公司，贷款算谁的？

这里贷款指的就是债权债务。（2011-03-09）

从我们卖学习工具的生意里好像也能体会到人口年龄层的变化。算起来房地产泡沫（如果有的话）大概到这些人需要买房子的时候就该差不多了。（2011-01-10）

\BookQuestion{网友：90后、00后人口数量锐减，是不是会对很多最终客户是年轻人，根基又不扎实的生意模式逐渐产生致命影响？}

会有影响，这个影响面会比大多数人想象的宽。（2015-05-23）

\BookQuestion{网友：对于今天的万科您怎么看？10年后的万科您怎么看？}

我本身对房地产不熟，很久没看过万科和房地产这个行业了。目前中国人口老龄化问题似乎在加重中，不知道对房地产未来有什么影响。（2019-03-14）

\BookQuestion{网友：大道曾经说过万科是个不错的公司，也建议过您的朋友等万科便宜一些再买。现在万科市盈率8.8，您还建议您的朋友买吗？}

那是快20年前说的了，那时候万科的债务还不是那么厉害。后来发现其债务很高后就没再关注过。（2024-03-31）
